\subsection{GELU and sigmoid-gated rows}
\label{sec:practical-feedforward}

The additive composition argument applies to updates other than tanh.
The following elementary constructions record the constants for two
common feed-forward choices. They use fresh randomness for each factor
in a product. Every closed gate returns a fresh fair sign.

Suppose normalized-stream coordinates satisfy $|a_i|\le G$, and fresh
signs of means $a_i/G$ are available. Let $u$ and $v$ be affine rows
with biases of magnitude at most one and row norms at most $s$.
Put $A=1+sG$ and $\rho=sG/2$. A sign of mean $u/A$ or $v/A$ uses
at most one normalized-stream request, by a signed row mixture.
Write
\[
 \sigma(v)=\frac1{1+e^{-v}},\qquad
 \operatorname{GELU}(v)=\frac v2
       \left(1+\operatorname{erf}(v/\sqrt2)\right).
\]

\begin{lemma}[Exact GELU and gated-product interfaces]
\label{lem:gelu-swiglu}
The following updates have exact normalized sign factories with the
stated output radii and expected normalized-stream request counts:
\[
\begin{array}{ccc}
 \toprule
 \textup{update}&\textup{radius}&\textup{expected requests, at most}\\ \midrule
 \operatorname{GELU}(v)&A&3/2+A^2/2\\
 u\operatorname{GELU}(v)&A^2&5/2+A^2/2\\
 uv\sigma(v)&A^2&2+2\rho(1+\rho)\\
 \bottomrule
\end{array}
\]
Their expected auxiliary bit costs are respectively
$O((1+A^2)B^4)$, $O((1+A^2)B^4)$, and
$O((1+\rho)^2B^4)$. All constants are absolute. A following affine
projection with bias bound $b_0$ and row norm $s_0$ changes a radius
$U$ to $b_0+s_0U$ and needs at most one such update request per
normalized projected sign.
\end{lemma}
\paragraph{Proof idea.}
Write GELU as a linear term plus a linear term times an error
function. Independent signs realize these products exactly. For SwiGLU,
a logistic gate selects a product of two fresh affine-row signs; the final
projection is a signed mixture.

\begin{proof}
Let $z=v/A$. Lemma~\ref{lem:erffactory}, with $c=A/\sqrt2$, gives
an erf sign of mean $\operatorname{erf}(v/\sqrt2)$ and source count
at most $1+A^2$. Each of those sources is one fresh $v/A$ row sign.
With probability one half, return a $v/A$ sign. With probability
one half, return the product of an independent $v/A$ sign and the
erf sign. Its mean is
\[
 \frac z2+\frac z2\operatorname{erf}(v/\sqrt2)
 =\frac{\operatorname{GELU}(v)}A.
\]
The expected source count is $1+(1+A^2)/2$. The finite-bit erf
implementation uses one Gaussian and a Poisson mean at most $A^2/2$;
the rational-plus-residual Poisson construction in
Lemma~\ref{lem:gaussianrmsbits} gives the displayed auxiliary bound.
Multiply this output by a fresh $u/A$ sign for the gated GELU row.

For the last row, the biased tanh factory produces a sign of mean
$\tanh(v/2)$ with at most $2\rho(1+\rho)$ normalized-stream requests.
Its positive outcome is a gate of probability
$(1+\tanh(v/2))/2=\sigma(v)$. On an open gate, multiply independent
signs of means $u/A$ and $v/A$. On a closed gate, return a fair sign.
The mean is $uv\sigma(v)/A^2$. Its two additional source requests
and the tanh factory's local bound prove the claimed costs.

Finally, a signed mixture over projection-row coefficients and a known
bias sign gives the normalized affine projection, with a fair-sign
completion if its supplied radius exceeds the exact row bound.
Every mixture is finite, every component factory terminates almost
surely, and each product uses independent fresh calls. Conditional
expectations prove exactness for every fixed input stream.
\end{proof}

Under a relative energy certificate, $G$ is a fixed normalization
radius, so all entries of the table are constants at fixed row norms.
Theorem~\ref{thm:pn-abstract} therefore still yields a polynomial depth
bound, with the corresponding update radius and source count. The
presence of these activations does not itself certify the energy floor
or bound the learned matrix norms.

\subsubsection{Reducing a supplied radius}

\begin{lemma}[Amplifying a sign mean with a fixed margin]
\label{lem:bipolar-amplification}
Let $C\ge1$ be rational. From independent signs of unknown mean $z$
with the promise $|z|\le1/(2C)$, one can generate an exact sign of
mean $Cz$ using at most $722C(C+1)\le1444C^2$ source requests in
expectation and $O(C^2B^4)$ expected auxiliary bit work. The factory
terminates almost surely.
\end{lemma}
\paragraph{Proof idea.}
Two affine changes of Bernoulli bias turn sign amplification into
two linear factories with explicit slack. Their exact affine identities
determine the output bias, and conditional charging multiplies their
source budgets.

\begin{proof}
The case $C=1$ returns a source sign. Otherwise let
\[
 p=(1+z)/2,\quad a=(C-1)/(2C),\quad
 D=\frac1{1-a}=\frac{2C}{C+1},\quad E=\frac{C+1}{2}.
\]
Apply the linear factory to a $(1-p)$-coin with multiplier $D$.
The promise gives
\[
 D(1-p)\le1-\frac1{2(C+1)}.
\]
Complement the result. Its bias is
$r=1-D(1-p)=(p-a)/(1-a)$. A second linear factory with multiplier $E$
gives bias
\[
 Er=C(p-a)=\frac{1+Cz}{2}\in[1/4,3/4].
\]
Read this bit as a sign. Huber's $9.5$ bound gives at most $38C$ source
requests for an $r$-coin, and at most $19(C+1)$ $r$-coins for the
second factory. Predictable charging yields their product,
$722C(C+1)$. The first factory has expected auxiliary work $O(CB^4)$;
the second has $O(CB^4)$. Charging the expected $O(C)$ first-stage
invocations proves the total $O(C^2B^4)$ bound. Exactness follows from
the displayed affine identities. Almost-sure termination follows from
the two linear factories and their almost surely finite call counts.
\end{proof}

\begin{corollary}[Reducing a supplied normalization radius]
\label{cor:radius-reduction}
Suppose a source returns signs of mean $F/R$, where a separately
certified bound gives $|F|\le U$. There is an exact sign sampler for
$F/(2U)$ with at most
\[
 \max\{1,361(R/U)^2\}
\]
expected source invocations and $O((1+(R/U)^2)B^4)$ auxiliary bit
work. Here $R,U>0$ are finitely encoded rationals.
\end{corollary}
\paragraph{Proof idea.}
A supplied radius can be increased by thinning. When it must be
decreased, the separate magnitude certificate provides the fixed margin
required by the sign-amplification lemma.

\begin{proof}
If $R\le2U$, use a known thinning coin of probability $R/(2U)$ and
return a fair sign on the closed branch. If $R>2U$, apply
Lemma~\ref{lem:bipolar-amplification} with $C=R/(2U)>1$ and
$z=F/R$. Its premise holds since $|F|/R\le U/R=1/(2C)$. The input-count
bound $1444C^2=361(R/U)^2$ proves the claim.
\end{proof}

\subsubsection{Depth lower bounds for GELU and SwiGLU}

We now carry the relative-floor lower construction of
Theorem~\ref{thm:pn-scaled-lower} over to GELU and SwiGLU blocks.
Lemma~\ref{lem:gelu-swiglu} and Corollary~\ref{cor:radius-reduction}
supply exact update factories at the declared radii.
Corollary~\ref{cor:real-ffn-scheduled} gives the GELU and SwiGLU case of
\eqref{eq:main-prenorm-lower} for every dyadic step schedule.

\begin{theorem}[A depth obstruction with GELU or SwiGLU]
\label{thm:real-ffn-lower}
Fix a dyadic $s\ge2$, and set
\[
 \delta=s-1,\qquad C_* =\frac{s}{4(s-1)},\qquad
 \gamma_* =\frac1{192s^2}.
\]
There is a positive constant $c_s^{\rm ff}$, depending only on $s$,
with the following property.
For every $n\ge4$, dyadic $0<\alpha\le1$, and $D$ with $\alpha D\ge4$,
there is a one-position pre-norm additive network of width $n$ and
depth $D$ whose feed-forward activation is either exact GELU or
SwiGLU, with all affine row norms at most $s$, biases at most one,
zero attention, fixed stabilizer $3$, and step $\alpha$ on every
sublayer, for which
\[
 \mathbb E Q\ge
 c_s^{\rm ff}
 \min\{(1+\alpha D)^{2\delta},n\}.
\]
An explicit value for the constant is given in the proof.
Preprocessing independent of the fresh query is arbitrary. The output
is a fixed final $\operatorname{RMSNorm}_3$ followed by logits
$\pm\mathcal N_3(h_D)_1/4$. These logits are bounded by one on the
entire input sign cube. Every normalization has the relative floor
$v_j\ge\gamma_*R_j^2$ for the valid supplied radii
$R_j=1+4s\alpha j$. Exact update factories at these radii have costs
depending only on $s$.
\end{theorem}
\paragraph{Proof idea.}
Both activations satisfy an exact antisymmetric identity that
recovers a linear signal. A constant hidden unit drives the carriers.
This embeds the normalized-gain lower construction into each feed-forward
architecture while preserving its full-cube energy floor.

\begin{proof}
Let $N$ be the largest power of two at most $n/2$, and put
$\eta=N/n\in(1/4,1/2]$. The first $N$ coordinates are features and
the remaining coordinates are carriers. Write
$\phi(v)=\operatorname{GELU}(v)$ or $\phi(v)=v\sigma(v)$.
Both choices satisfy
\[
 \phi(v)-\phi(-v)=v,\qquad \tfrac12<\phi(1)<1.
\]
The first identity follows from the oddness of the error function for
GELU, and from $\sigma(v)+\sigma(-v)=1$ for SiLU. Use three hidden
units with preactivations
\[
  \frac{s}{N}\sum_{i\le N}\mathcal N_3(h)_i,\qquad
  -\frac{s}{N}\sum_{i\le N}\mathcal N_3(h)_i,\qquad 1.
\]
Every feature projection row is $(1,-1,0)$ and every carrier row is
$(0,0,-1)$. Extra hidden coordinates may be zero. In the SwiGLU case
$u\,v\,\sigma(v)$, the displayed preactivations are the gate arguments
$v$, and the ungated factor $u$ of each of these three hidden units is
the constant one. Thus the SwiGLU units equal $\phi$ exactly. The
input rows have norms at most $s$, and the projection rows have norms
at most two, which is allowed because $s\ge2$.

For the lower-bound distribution let the features be independent
unknown signs $x_i$ and set every initial carrier to $-1$. Put
$M=N^{-1}\sum_i x_i$ and $c=\phi(1)$. After $j$ blocks the feature
coordinates are $x_i-M+F_j(M)$ and all carriers are
$C_j=-(1+\alpha jc)$, where $F_0(z)=z$ and
\[
 F_{j+1}(z)=F_j(z)+\alpha
 \frac{sF_j(z)}{
 \sqrt{3+\eta(1-z^2+F_j(z)^2)+(1-\eta)C_j^2}}.
\]
This follows from the exact difference identity above. These
recurrences are odd in $z$, and $F_j(z)\ge z\ge0$ when $z\ge0$.

For every input, the mean of the feature coordinates has magnitude at
most $\sqrt{v_j/\eta}$, by Cauchy--Schwarz and the definition of the
RMS denominator. The feature update therefore has magnitude at most
$s/\sqrt\eta\le2s$, and the carrier update has magnitude below one.
In particular, $|h_{j,i}|\le1+2s\alpha j$. The coordinate update
bound $2s$ is valid on the entire sign cube.

Put $r_j=1+\alpha j$. For arbitrary initial carrier signs their values
are $e_i-\alpha jc$, with $c>1/2$. The full-cube calculation in
Corollary~\ref{cor:pn-scaled-padding}, using the carrier fraction
$1-\eta\ge1/2$ and the stabilizer $3$, gives
\[
 v_j\ge r_j^2/12.
\]
Because $R_j=1+4s\alpha j\le4sr_j$, this implies the stated
$v_j\ge R_j^2/(192s^2)$. To obtain the exact update factory at radius
$4s$, first use Lemma~\ref{lem:gelu-swiglu} and its affine projection
to obtain a normalized sign with some fixed radius depending only
on $s$ and the supplied floor $\gamma_*$. The actual update bound is
$2s$. Corollary~\ref{cor:radius-reduction} supplies an exact sign at
radius $4s$, still with costs depending only on $s$. This accounts
for the declared radii without requiring the generic activation
bound to equal the true update range.

For the growth comparison, normalize the feature mean by $r_j$,
writing $u_j=F_j/r_j$ and $a_j=\alpha/r_{j+1}$. After the burn-in
$j_0=\lceil3/\alpha\rceil$, one has $4\le r_{j_0}\le5$ and
$u_{j_0}(z)\ge z/5$. The nonsignal part of the normalized denominator
satisfies
\[
 b_j(z)\le\frac{7}{2r_j^2}+\frac34\le\frac{31}{32}<1.
\]
Consequently, for $z\ge0$,
\[
 u_{j+1}\ge(1-a_j)u_j+
          a_j\frac{s u_j}{\sqrt{1+u_j^2/2}}.
\]
Apply the convexity comparison used in
Theorem~\ref{thm:pn-scaled-lower}, now with
$g_j=1+\delta a_j$ and $d_j=a_js/(2g_j)$. Its reciprocal-square
sum has coefficient $C_*=s/(4\delta)$. Thus, for
$G=\prod_{j=j_0}^{D-1}g_j$,
\[
 u_D(z)\ge
 \frac{G(z/5)}{\sqrt{1+C_*(G^2-1)(z/5)^2}},\qquad
 G\ge e^{-\delta^2/8}\left(\frac{1+\alpha D}{6}\right)^\delta.
\]
The latter product bound is uniform in $\alpha$.

It remains to control the fixed output readout. On every sign input
the feature decomposition $h_i=x_i-M+F$ still holds, even when the
initial carrier signs vary. Hence $|h_i|\le |F|+2$ and
$v\ge3+\eta F^2\ge3+F^2/4$. The readout calculation in the proof of
Theorem~\ref{thm:pn-normalized-gain} then shows that the absolute logit
is at most $1/\sqrt3<1$.
On the lower-bound subdomain, conditional on $M=z$, set
$L(z)=4\sqrt{v_D(z)}$. The coordinate bound gives
\[
 L(z)\le4\sqrt{3+(2s r_D)^2}\le8(s+1)r_D.
\]
The conditional shift proof in
Corollary~\ref{cor:pn-scaled-finalnorm-lower} applies verbatim:
the actual and intermediate head arguments lie in $[-1,1]$,
$\operatorname{sech}^2(1)>1/4$, and
\[
 \overline g(z)\ge\frac{F_D(z)}{4L(z)}
 \ge\frac{u_D(z)}{32(s+1)}\qquad(z\ge0).
\]
Weighted Jensen and the Rademacher fourth-moment bound, as in the
proof of Theorem~\ref{thm:pn-scaled-lower}, now give
\begin{gather*}
 H'(0)\ge
 \frac{G}{160(s+1)\sqrt{1+3C_*G^2/(25N)}},\\
 \mathbb E Q\ge
 \frac{\min\{G^2,N\}}{25600(s+1)^2(1+3C_*/25)}.
\end{gather*}
Substitute the bound for $G$ and use $N>n/4$. This proves the theorem
with
\[
 c_s^{\rm ff}=\frac{e^{-\delta^2/4}6^{-2\delta}}
                     {102400(s+1)^2(1+3C_*/25)}.
\]
\end{proof}

\begin{corollary}[GELU and SwiGLU with a prescribed step schedule]
\label{cor:real-ffn-scheduled}
Theorem~\ref{thm:real-ffn-lower} also holds for every depth $D\ge0$
and every prescribed dyadic schedule $a_j\in[0,1]$, with $\alpha D$
replaced by $\tau_D=\sum_{j<D}a_j$ and the condition $\alpha D\ge4$
dropped, the supplied radius replaced by $R_j=1+4s\sum_{k<j}a_k$, and
a positive lower-bound prefactor depending only on $s$. Its value is
given in the proof.
\end{corollary}
\paragraph{Proof idea.}
The activation identities and output-radius bounds hold at every
step size. Replacing elapsed depth by cumulative carrier displacement
therefore transfers the scheduled normalized-gain argument directly.
A small total step needs only the head bound.

\begin{proof}
The states and relative floors depend on the step schedule through
the cumulative carrier displacement and the feature recurrence.
Put $\tau_j=\sum_{k<j}a_k$ and $r_j=1+\tau_j$. If $\tau_D\ge4$, the
proof of Corollary~\ref{cor:pn-scheduled-lower} applies with the
same linear-update comparison constant $C_*$. The head and source
radius arguments in Theorem~\ref{thm:real-ffn-lower} are unchanged.
The resulting prefactor is
\[
 \frac{e^{-\delta/2-\delta^2/4}5^{-2\delta}}
 {102400(s+1)^2(1+3C_*/25)}.
\]

If $\tau_D<4$, no growth comparison is needed. The head bound in
Theorem~\ref{thm:real-ffn-lower}, with $L(z)\le8(s+1)r_D$ and
$F_D(z)\ge z$, gives $\overline g(z)\ge z/(32(s+1)r_D)$ for $z\ge0$.
Hence $H'(0)\ge1/(32(s+1)r_D)$, and since $r_D<5$,
\[
 \mathbb EQ\ge H'(0)^2>\frac1{25600(s+1)^2}.
\]
Because $\min\{(1+\tau_D)^{2\delta},n\}<5^{2\delta}$, the displayed
prefactor also covers this case.
\end{proof}
